\documentclass[11pt,a4paper]{article}
\usepackage[czech]{babel}
\usepackage[utf8]{inputenc}
\usepackage{times}
\usepackage{url}
\usepackage[textwidth=15.2cm,textheight=23cm]{geometry}
\usepackage{xcolor}

\usepackage{graphicx}

%\usepackage{fancyvrb}
%\DefineVerbatimEnvironment{verbatim}{Verbatim}{}

\usepackage[bf]{caption}

\usepackage[hyperindex,
  plainpages=false,
  pdftex,
  colorlinks,
  pdfborder={0 0 0},
  pdfpagelabels]{hyperref}

\pdfcompresslevel=9

\newcommand{\myincludegraphics}[4]{
  \begin{figure}[!h]
  \centering
  \includegraphics[#1]{#2}
  \caption{#3.} \label{#4}
  \end{figure}
}

% titulní stránka a obsah
\newcommand{\titlepageandcontents}{
  % credits for template go to: Martin Striz
\begin{titlepage}

\vspace*{1cm}

\begin{figure}
  \centering
  \includegraphics[height=6cm]{images/fit.pdf}
\end{figure}

\vspace*{5mm}

\begin{center}
\begin{Large}
Projekt do předmětu GMU -- Grafické a multimediální procesory
\end{Large}
\end{center}

\vspace*{5mm}

\begin{center}
\begin{Huge}
Neuronové sítě v OpenCL
\end{Huge}
\end{center}

\vspace*{1cm}

\begin{center}
\begin{Large}
\today
\end{Large}
\end{center}

\vfill

\begin{flushleft}
\begin{large}
\begin{tabular}{ll}

\bf Řešitel:\hspace{3mm} & Jakub Zárybnický (\verb_xzaryb00@stud.fit.vutbr.cz_) \\
& Fakulta Informačních Technologií \\
& Vysoké Učení Technické v~Brně

\end{tabular}
\end{large}
\end{flushleft}

\end{titlepage}

% vim:set ft=tex expandtab enc=utf8:


  \pagestyle{plain}
  \pagenumbering{roman}
  \setcounter{page}{1}
  %\tableofcontents

  \newpage
  \pagestyle{plain}
  \pagenumbering{arabic}
  \setcounter{page}{1}
}

\def\uv#1{\iflanguage{english}{``#1''}%
                              {\quotedblbase #1\textquotedblleft}}%

% vim:set ft=tex expandtab enc=utf8:

\usepackage{hyperref}
\begin{document}
\titlepageandcontents

%---------------------------------------------------------------------------
\section{Zadání}

Formální zadání:
\begin{itemize}
  \item Nastudujte neuronové sítě (back propagation, inference, ...)
  \item Implementujte framework umožňující konfiguraci topologie sítě, nastavování vlastních aktivačních funkcí, trénování).
  \item Implementujte jednoduchou aplikaci využívající framework.
\end{itemize}

Interpretace:

\begin{itemize}
\item Několik pokusů o řešení: CUDA (následné zjištění, že se CUDA nedá
  emulovat), Vulkan (pomocí compute shaders), raw OpenCL v C++ (po správném
  přečtení zadání ``v OpenCL''), raw OpenCL v Haskellu a poslední varianta:
\item Poté, co se mé řešení začalo nebezpečně podobat jednomu existujícímu
  frameworku jsem zvolil alternativu takovou, že tento framework využiji a
  vylepším/paralelizuji.
\end{itemize}

\section{Použité technologie}

\begin{itemize}
  \item OpenCL
  \item HMatrix
  \item Grenade
    \\
  \item Namísto CMake použit Nix, který narozdíl od CMake dobře podporuje Haskell
  \item Žádné další závoslosti,
\end{itemize}

\section{Použité zdroje}

\begin{itemize}
  \item \url{https://github.com/HuwCampbell/grenade}
  \item \url{https://github.com/IFCA/opencl}
  \item \url{https://github.com/acowley/CLUtil}
\end{itemize}

%---------------------------------------------------------------------------
\section{Způsob akcelerace}

\begin{itemize}
  \item Paralelizace procesu trénování i inference sítě v OpenCL (matrix-vector
    multiplication, matrix-matrix multiplication, outer vector product)
  \item Konkrétně implementace tří maticových algoritmů pro FCNN
  \item Následné zjištění, že soupeřím s BLAS/LAPACK v původní verzi software a
    ladění výkonu FFI mezi Haskell a OpenCL
\end{itemize}

První pokusy o zrychlení pouze \texttt{cbits} částí frameworku --- konkrétně
gradient descent a konvoluční algoritmy --- vedly k pětinásobnému zpomalení a
rychle jsem je zavrhl - \textit{overhead} neustálého kopírování dat mezi
programem a OpenCL byl příliš vysoký.

Implementoval jsem namísto toho typ vrstvy neuronů, která celá věží na OpenCL,
ale je zároveň interoperabilní se zbytkem kódu psaného čistě v Haskellu.

\section{Ovládání vytvořeného programu}

Je potřeba nainstalovat Nix, konkrétně \texttt{nixFlakes}, experimentální Nix
verze 3.0, který pak nainstaluje veškeré potřebné závislosti.

Jelikož se má jednat o framework, demonstrační program samotný je triviální a dá
se spustit v kořenovém adresáři uvnitř prostředí \texttt{nix shell} pomocí

\begin{verbatim}
cabal run grenade-examples:exe:feedforward
\end{verbatim}

popř. pro profilování je možné použít

\begin{verbatim}
cabal run --enable-profiling -- grenade-examples:exe:feedforward +RTS -p -s
\end{verbatim}

\section{Vyhodnocení}

Testováno na \texttt{intel-ocl} na vestavěné grafické kartě, jelikož nemám
přístup k GPU zařízením.

Rychlost akcelerovaných vrstev je přibližně konstatní při rostoucí velikosti,
nebo roste zanedbatelně, zatímco rychlost vrstev používajících BLAS roste
lineárně --- zřejmě FFI \textit{overhead} je při použití příliš malých matic
příliš vysoký, protože maticové operace pouze pomocí BLAS bez nutnosti
kopírování matic jsou znatelně rychlejší. Při velikosti vrstev kolem 700 neuronů
už se rychlosti vyrovnaly.

%---------------------------------------------------------------------------
\section{Rozdělení práce v týmu}

Sólový projekt.

\end{document}
% vim:set ft=tex expandtab enc=utf8:
%%% Local Variables:
%%% mode: latex
%%% TeX-master: t
%%% End:
